\documentclass[11pt,a4paper]{article}
\usepackage[margin=1in]{geometry}
\usepackage{amsmath,amssymb,amsthm}
\usepackage{graphicx}
\usepackage{xcolor}
\usepackage{booktabs}
\usepackage{hyperref}
\usepackage{enumitem}
\usepackage{fancyhdr}
\usepackage{tcolorbox}
\usepackage{titlesec}
\usepackage{microtype}

% Colors
\definecolor{ctiblue}{RGB}{31,78,121}
\definecolor{ctiorange}{RGB}{196,89,17}
\definecolor{ctigreen}{RGB}{39,124,52}
\definecolor{ctired}{RGB}{180,30,30}
\definecolor{ctipurple}{RGB}{100,50,150}
\definecolor{lightblue}{RGB}{230,240,250}
\definecolor{lightorange}{RGB}{255,245,230}
\definecolor{lightgreen}{RGB}{230,250,235}
\definecolor{lightred}{RGB}{255,235,235}

% Hyperlinks
\hypersetup{colorlinks=true, linkcolor=ctiblue, urlcolor=ctiorange, citecolor=ctiblue}

% Header/footer
\pagestyle{fancy}
\fancyhf{}
\fancyhead[L]{\small\textcolor{ctiblue}{CTI Universal Law --- First-Principles Guide}}
\fancyhead[R]{\small\textcolor{ctiblue}{\thepage}}
\renewcommand{\headrulewidth}{0.4pt}

% Section formatting
\titleformat{\section}{\Large\bfseries\color{ctiblue}}{}{0em}{}[\titlerule]
\titleformat{\subsection}{\large\bfseries\color{ctiorange}}{}{0em}{}
\titleformat{\subsubsection}{\normalsize\bfseries\color{ctigreen}}{}{0em}{}

% Boxes
\newtcolorbox{keybox}[1][]{colback=lightblue, colframe=ctiblue, title={\textbf{#1}}, fonttitle=\bfseries\color{white}, coltitle=white, colbacktitle=ctiblue, sharp corners, boxrule=0.5pt}
\newtcolorbox{derivbox}[1][]{colback=lightorange, colframe=ctiorange, title={\textbf{#1}}, fonttitle=\bfseries\color{white}, coltitle=white, colbacktitle=ctiorange, sharp corners, boxrule=0.5pt}
\newtcolorbox{resultbox}[1][]{colback=lightgreen, colframe=ctigreen, title={\textbf{#1}}, fonttitle=\bfseries\color{white}, coltitle=white, colbacktitle=ctigreen, sharp corners, boxrule=0.5pt}
\newtcolorbox{warnbox}[1][]{colback=lightred, colframe=ctired, title={\textbf{#1}}, fonttitle=\bfseries\color{white}, coltitle=white, colbacktitle=ctired, sharp corners, boxrule=0.5pt}

\DeclareMathOperator{\logit}{logit}

\begin{document}

% Title page
\begin{titlepage}
\centering
\vspace*{2cm}
{\Huge\bfseries\color{ctiblue} The CTI Universal Law}\\[0.8cm]
{\LARGE\color{ctiorange} A First-Principles Study Guide}\\[1.5cm]
{\large Everything you need to understand the math,\\the evidence, and the open directions.}\\[3cm]

\begin{tcolorbox}[colback=lightblue, colframe=ctiblue, width=0.85\textwidth, sharp corners]
\centering\large
There is a single equation that tells you how well \textbf{any} classifier\\
(neural network, biological brain, anything) will perform,\\
using only the \textbf{geometry} of how it organizes information ---\\
and it's derivable from first principles.\\[0.5cm]
\Large\bfseries\color{ctiblue}
$\displaystyle\logit(q_{\mathrm{norm}}) = \alpha\,\kappa_{\mathrm{nearest}} - \beta\log(K{-}1) + C$
\end{tcolorbox}

\vfill
{\large Devansh $\cdot$ March 2026}
\end{titlepage}

\tableofcontents
\newpage

% ===================================================================
\section{The Setup: What Are We Measuring?}
% ===================================================================

\subsection{The fundamental question}

Given a trained neural network (or a biological brain), how good are its internal representations at distinguishing between categories?

We don't care about the final classifier head. We care about the \textbf{geometry of the embeddings} --- the spatial layout of how the model organizes data points internally. If two classes are well-separated in embedding space, the model can tell them apart easily. If they're tangled up, it can't.

\subsection{The probe: 1-Nearest Neighbor (1-NN)}

To measure representation quality without adding any learnable parameters, we use the simplest possible classifier: \textbf{1-nearest neighbor}.

\begin{keybox}[1-NN Classification]
Given a test point $\mathbf{x}$:
\begin{enumerate}[nosep]
\item Find its nearest neighbor in the training set (by Euclidean distance in embedding space)
\item Predict $\mathbf{x}$ has the same label as that neighbor
\end{enumerate}
1-NN accuracy depends \textbf{only} on the geometry of the embeddings. No probe capacity. No optimization. Pure geometric signal.
\end{keybox}

\textbf{Why not a linear probe?} A linear probe mixes geometry with the probe's own optimization. It can mask bad geometry if the probe has enough capacity, or fail with good geometry if the hyperparameters are wrong. 1-NN removes this confounder entirely: the accuracy \emph{is} the geometry.

% ===================================================================
\section{The Key Quantities}
% ===================================================================

\subsection{$q_{\mathrm{norm}}$: Chance-Normalized Accuracy}

Raw 1-NN accuracy is misleading across different $K$. With $K{=}2$, random guessing gives 50\%. With $K{=}100$, random gives 1\%. We normalize:

\begin{keybox}[Chance-Normalized Accuracy]
\vspace{-0.3cm}
\begin{equation}
q_{\mathrm{norm}} = \frac{\mathrm{acc}_{1\mathrm{NN}} - 1/K}{1 - 1/K}
\end{equation}
\begin{itemize}[nosep]
\item $q_{\mathrm{norm}} = 0$: exactly at chance (no geometric signal)
\item $q_{\mathrm{norm}} = 1$: perfect classification
\item $q_{\mathrm{norm}} < 0$: below chance (rare, adversarial geometry)
\end{itemize}
\end{keybox}

\textbf{Example}: If $K{=}10$ and 1-NN accuracy $= 55\%$:
\[
q_{\mathrm{norm}} = \frac{0.55 - 0.10}{1 - 0.10} = \frac{0.45}{0.90} = 0.50
\]

\subsection{$\kappa_{\mathrm{nearest}}$: The Geometric Signal-to-Noise Ratio}

This is the single most important number in the entire theory.

\begin{keybox}[$\kappa_{\mathrm{nearest}}$ --- The Key Quantity]
\vspace{-0.3cm}
\begin{equation}
\kappa_{\mathrm{nearest}}(k) = \min_{j \neq k} \frac{\|\boldsymbol{\mu}_j - \boldsymbol{\mu}_k\|}{\sigma_W \sqrt{d}}
\end{equation}
\textbf{Numerator}: distance to the \emph{nearest} other class centroid (the hardest competitor)\\
\textbf{Denominator}: expected within-class L2 radius (isotropic null)\\
\textbf{Result}: dimensionless signal-to-noise ratio for class separation
\end{keybox}

Where:
\begin{itemize}[nosep]
\item $\boldsymbol{\mu}_k$ = centroid (mean embedding) of class $k$
\item $\min_{j \neq k}$ = take the \textbf{nearest} other class (hardest competitor)
\item $\sigma_W$ = pooled within-class per-dimension standard deviation
\item $d$ = embedding dimension
\end{itemize}

\subsubsection{Why ``nearest''?}

Because in a 1-NN race, the class you're most likely to confuse a point with is the \textbf{nearest other class}. Far-away classes are irrelevant. This is the \emph{bottleneck} --- the weakest link.

\subsubsection{Intuition}

Think of $\kappa_{\mathrm{nearest}}$ as a ``contrast ratio'':
\begin{itemize}[nosep]
\item $\kappa \ll 1$: noise swamps signal $\to$ classification at chance
\item $\kappa \sim 1$: signal and noise comparable $\to$ transition zone
\item $\kappa \gg 1$: signal dominates $\to$ near-perfect classification
\end{itemize}

\subsubsection{Concrete example}

Pythia-160M on AGNews ($K{=}4$), final layer: centroid distances between news categories vary by pair. The nearest pair: $\|\boldsymbol{\mu}_\mathrm{sports} - \boldsymbol{\mu}_\mathrm{world}\| = 15.2$. With $\sigma_W\sqrt{768} = 14.3$: $\kappa_\mathrm{nearest} = 15.2/14.3 = 1.06$.

\subsection{$K$: Number of Classes}

$K$ enters through $\log(K{-}1)$ because more classes = more competitors. Doubling classes doesn't double difficulty --- it increases it \emph{logarithmically} (from the Gumbel extreme-value structure).

% ===================================================================
\section{Why Logit? Why Linear?}
% ===================================================================

\subsection{The logit transformation}

The logit maps probability $p \in (0,1)$ to $\mathbb{R}$:
\begin{align}
\logit(p) &= \log\frac{p}{1-p} \\
\sigma(x) &= \frac{1}{1+e^{-x}} \quad \text{(inverse: sigmoid)}
\end{align}

\begin{derivbox}[The Law in Two Forms]
\textbf{Logit form:}
\[
\logit(q_\mathrm{norm}) = \alpha\,\kappa_\mathrm{nearest} - \beta\log(K{-}1) + C
\]
\textbf{Equivalent sigmoid form:}
\[
q_\mathrm{norm} = \frac{1}{1 + \exp\!\bigl(-(\alpha\,\kappa_\mathrm{nearest} - \beta\log(K{-}1) + C)\bigr)}
\]
\end{derivbox}

\textbf{This is NOT a logistic regression.} It's a structural prediction from extreme value theory. The sigmoid appears because of the Gumbel race probability structure.

% ===================================================================
\section{The Gumbel Race: Where the Law Comes From}
% ===================================================================

This is the core derivation. Seven steps from first principles.

\subsection{Step 1: 1-NN as a race}

When classifying query $\mathbf{x}$ from class $k$ using 1-NN:

\begin{derivbox}[The 1-NN Race]
Define:
\begin{align}
M_+ &= \min_{i:\,y_i=k,\,i\neq x} \|\mathbf{x} - \mathbf{x}_i\| \quad \text{(nearest same-class)} \\
M_j &= \min_{i:\,y_i=j} \|\mathbf{x} - \mathbf{x}_i\| \quad \text{(nearest from class $j$, for $j\neq k$)}
\end{align}
\textbf{1-NN is correct} $\iff$ $M_+ < \min_{j \neq k} M_j$

The correct class must ``win the race'' by having the smallest minimum distance.
\end{derivbox}

\subsection{Step 2: Extreme Value Theory (EVT)}

The minimum of $n$ i.i.d.\ random variables, after centering and scaling, converges to one of exactly three universal limit distributions (Fisher--Tippett theorem, 1928):

\begin{enumerate}[nosep]
\item \textbf{Gumbel} --- exponential-type tails (Gaussian, etc.) $\leftarrow$ \textbf{our case}
\item \textbf{Fr\'echet} --- heavy tails
\item \textbf{Weibull} --- finite support
\end{enumerate}

For Gaussian/sub-Gaussian class conditionals (which neural embeddings approximately satisfy), minimum distances are in the \textbf{Gumbel domain of attraction}.

\begin{keybox}[The Gumbel Distribution]
\vspace{-0.3cm}
\[
P(G \leq x) = \exp\!\bigl(-\exp(-(x-\mu)/s)\bigr)
\]
$\mu$ = location (decreases as $O(\sqrt{\log n})$ with more samples)\\
$s$ = scale (stays roughly constant)
\end{keybox}

\subsection{Step 3: Setting up the Gumbel race}

After EVT normalization:
\begin{align}
M_+ &\sim \mathrm{Gumbel}(\mu_+, s) \quad \text{(same-class nearest)} \\
M_j &\sim \mathrm{Gumbel}(\mu_j^{(-)}, s) \quad \text{(class-$j$ nearest)}
\end{align}

\textbf{Key}: all Gumbels share approximately the \textbf{same scale} $s$ (within-class spread is similar). They differ only in \textbf{locations} (how far each class centroid is).

\subsection{Step 4: The Gumbel Race Identity}

\begin{derivbox}[The Magical Property of Gumbel Distributions]
If $G_0, G_1, \ldots, G_{K-1}$ are independent Gumbels with common scale $s$ and locations $\mu_0, \mu_1, \ldots, \mu_{K-1}$:
\[
\boxed{P(G_0 < \min_{j \geq 1} G_j) = \frac{1}{1 + \sum_{j \geq 1} \exp\!\bigl(-(\mu_j - \mu_0)/s\bigr)}}
\]
This is \textbf{exact}. Not approximate. The probability of one Gumbel beating all others is a \textbf{logistic function} of the location gaps.
\end{derivbox}

\textbf{Why this is extraordinary}: for most distributions, computing $P(\min \text{ of } K \text{ things})$ requires intractable $K$-dimensional integrals. For Gumbels, it factorizes into a simple closed-form sigmoid.

\subsection{Step 5: Apply to 1-NN}

Set $G_0 = M_+$ and $G_j = M_j$:
\[
q = P(\text{correct}) = \frac{1}{1 + \sum_{j \neq k} \exp(-\Delta_{k,j}/s)}
\]
where $\Delta_{k,j} = \mu_j^{(-)} - \mu_+$ is the ``gap'' for each competitor.

\subsection{Step 6: Nearest-class dominance + sparse hierarchy}

\begin{derivbox}[Two Competition Regimes]
\textbf{Dense (ETF/equidistant):} All $K{-}1$ competitors equidistant $\to$ $\sum = (K{-}1)e^{-\Delta/s}$ $\to$ $\beta = 1$

\textbf{Hierarchical (real world):} Only $\sqrt{K}$ effective competitors active $\to$ $\beta = 0.5$

\textbf{Empirical:} $\hat\beta = 0.478 \approx 0.5$ confirms \textbf{sparse competition}
\end{derivbox}

\subsection{Step 7: Linear gap-to-kappa map}

Under isotropic geometry, the gap is approximately linear in $\kappa$:
\[
\Delta_{\mathrm{nearest}} / s \approx a_1 \kappa_\mathrm{nearest} + a_0
\]

Combining everything:
\begin{resultbox}[The CTI Universal Law --- Derived]
\vspace{-0.3cm}
\[
\logit(q_\mathrm{norm}) = \underbrace{\frac{a_1}{s}}_{\alpha} \kappa_\mathrm{nearest} - \beta\log(K{-}1) + \underbrace{\frac{a_0}{s}}_{C}
\]

\textbf{Derived}: The functional form (logistic in $\kappa$, logarithmic in $K$)\\
\textbf{Estimated}: The constants $(\alpha, \beta, C)$
\end{resultbox}

% ===================================================================
\section{Neural Collapse: Why the Geometry Is Universal}
% ===================================================================

\subsection{What is Neural Collapse?}

Neural Collapse (NC) is a phenomenon discovered by Papyan, Han, and Donoho (PNAS 2020). At the terminal phase of training:

\begin{keybox}[The Four Properties of Neural Collapse]
\begin{enumerate}[nosep]
\item \textbf{NC1}: Within-class variability collapses (all points cluster tightly around centroid)
\item \textbf{NC2}: Class centroids form an Equiangular Tight Frame (ETF) --- maximally equidistant, like a regular simplex
\item \textbf{NC3}: Classifier weights align with centroids
\item \textbf{NC4}: Predictions converge to nearest-centroid classifier
\end{enumerate}
\end{keybox}

\subsection{Why NC matters for CTI}

NC provides the geometric \textbf{reason} why the law works across architectures:

\textbf{At exact NC}: all class pairs equidistant $\to$ $\kappa_\mathrm{nearest} = \kappa_\mathrm{spec}$, within-class isotropic $\to$ Gaussian assumption satisfied, symmetric Gumbel race $\to$ $\beta = 1$.

\textbf{Near NC} (real networks): centroids $\approx$ regular simplex, noise $\approx$ shared across classes, bottleneck pair drives classification.

\begin{keybox}[The Key Insight]
$\kappa_\mathrm{nearest}$ measures \textbf{how close} the representation is to Neural Collapse geometry. The law maps this geometric proximity to classification quality.

Different architectures converge to similar NC geometry on the same dataset $\to$ similar $\kappa$ $\to$ similar quality. The architecture determines the \emph{path}; the law governs the \emph{endpoint}.
\end{keybox}

% ===================================================================
\section{The Equicorrelation $\rho$: A Deeper Invariant}
% ===================================================================

\subsection{What $\rho$ measures}

$\rho$ is the average $\Sigma_W$-whitened cosine similarity between centroid-difference vectors:

\begin{enumerate}[nosep]
\item Take the within-class covariance $\Sigma_W$
\item Whiten: $\Sigma_W^{-1/2}(\boldsymbol{\mu}_j - \boldsymbol{\mu}_k)$ for all pairs
\item Compute pairwise cosine similarities of these whitened vectors
\item Average
\end{enumerate}

\begin{itemize}[nosep]
\item $\rho = 0$: orthogonal competition directions
\item $\rho = 1/2$: regular simplex (exact ETF)
\item $\rho = 1$: degenerate 1D competition
\end{itemize}

\subsection{The remarkable finding: $\rho \approx 0.46$ everywhere}

\begin{resultbox}[Cross-Modal Equicorrelation --- The Tightest Invariant]
\begin{center}
\begin{tabular}{lc}
\toprule
\textbf{Modality} & $\boldsymbol{\rho}$ \\
\midrule
NLP decoders (11 architectures) & 0.463 \\
Audio speech (WavLM, HuBERT) & 0.455--0.464 \\
Vision (ViT-Base, ResNet50) & 0.458--0.467 \\
Mouse visual cortex (Allen) & 0.466 \\
NLP encoders (BERT, ELECTRA) & 0.446--0.456 \\
\midrule
\textbf{Pooled (6 modalities)} & \textbf{0.462, CV = 1.0\%} \\
\bottomrule
\end{tabular}
\end{center}
Regular simplex prediction: $\rho = 0.500$. Observed: 8\% below perfect NC.
\end{resultbox}

\subsection{$\alpha$ prediction from $\rho$}

From the Gumbel race, the competition effective dimension is $d_\mathrm{eff,comp} = 1/(1-\rho)$, and:
\[
\alpha_\mathrm{pred} = \sqrt{\frac{4}{\pi}} \cdot \frac{1}{\sqrt{1-\rho}} = \frac{1.128}{\sqrt{0.54}} = 1.534
\]

\textbf{Observed}: $\alpha = 1.477$. \textbf{Error}: $+3.8\%$. \textbf{Zero free parameters.}

\begin{warnbox}[Caveat]
This only works for NLP decoders. For ViT ($\alpha \approx 4.5$) with the same $\rho \approx 0.46$, the prediction gives $\sim$1.5, not 4.5. A second renormalization factor (related to $\Sigma_W$ structure) amplifies $\alpha$ for continuous-signal modalities. This is an \textbf{open problem}.
\end{warnbox}

% ===================================================================
\section{The Three Constants: $\alpha$, $\beta$, $C$}
% ===================================================================

\subsection{$\alpha$: The slope (signal amplification)}

$\alpha$ converts geometric SNR into classification log-odds. Bigger $\alpha$ = steeper sigmoid.

\begin{center}
\begin{tabular}{lcc}
\toprule
\textbf{Signal Type} & $\boldsymbol{\alpha}$ & \textbf{Modalities} \\
\midrule
Discrete tokens (NLP decoders) & $\sim$1.48 & Pythia, Qwen, OLMo, \ldots \\
Continuous signals & $\sim$3.3--5.0 & ViT ($\sim$4.5), CNN ($\sim$4.0), Audio ($\sim$4.7) \\
NLP encoders & $\sim$7.7 & BERT, DeBERTa, ELECTRA \\
\bottomrule
\end{tabular}
\end{center}

\textbf{Within NLP decoders}: CV = 2.3\% across 12 architectures --- effectively a \textbf{true constant}. This CV is consistent with LOAO estimation noise (expected $\sim$2.8\%).

\subsection{$\beta$: The competition penalty}

\begin{keybox}[$\beta$ Values]
\textbf{Theory} (all $K{-}1$ competitors active, dense ETF): $\beta = 1.0$\\
\textbf{Empirical} (sparse hierarchy, $N_\mathrm{eff} \approx \sqrt{K}$): $\hat\beta = 0.478 \approx 0.5$

Only $\sim$$\sqrt{K}$ semantically relevant competitors are active. Pre-registered test confirms $\beta{=}0.5$ beats $\beta{=}1.0$ by 10.1\% LOAO MAE.
\end{keybox}

\subsection{$C$: The intercept (task difficulty)}

$C$ encodes base difficulty of a dataset. It is \textbf{NOT universal} --- each new dataset requires at least one calibration measurement. LODO (leave-one-dataset-out) mean $r = 0.125$: the intercept can't be predicted across tasks.

% ===================================================================
\section{Cross-Modal Validation}
% ===================================================================

\subsection{Vision}

\textbf{ViT-Large on CIFAR-10}: $R^2 = 0.964$, same logit-linear form. $\alpha_\mathrm{ViT} \approx 4.5$.

\textbf{ResNet50 on CIFAR-100}: $r = 0.749$, $\alpha_\mathrm{CNN} \approx 4.4$.

\subsection{Audio}

Seven frozen speech models from 4 architectures on Speech Commands ($K{=}36$):
$r = 0.898$, $p = 0.006$, $\alpha_\mathrm{audio} = 4.669$.

All continuous-signal modalities converge: $\alpha \approx 4$--$5$.

\subsection{Biological neural systems}

\textbf{Allen Neuropixels} ($K{=}118$, 32 mice): 30/32 sessions PASS ($r > 0.50$). Mean $r = 0.736$.

The law form is preserved in biological neurons. The constant differs ($A_\mathrm{bio}$ is $15$--$34\times$ smaller), but the geometric skeleton is the same.

\textbf{Cross-modal $\rho$}: 0.466 in cortex, within 1\% of NLP decoders. Near-simplex geometry is \textbf{substrate-independent}.

\begin{warnbox}[Where biology fails]
\textbf{Human fMRI} (NSD, $K{=}12$): 1-NN at chance, $r = 0.12$, $p = 0.18$. Raw BOLD voxels lack sufficient SNR. The law requires above-chance 1-NN as a prerequisite.
\end{warnbox}

% ===================================================================
\section{The Generation Law}
% ===================================================================

\subsection{The key insight}

Next-token prediction (NTP) is $V$-way classification: the model outputs logits $z_v = \mathbf{w}_v^\top \mathbf{h}(x)$ and the correct token wins the softmax ``race.'' The \textbf{same Gumbel race mechanism applies}.

\begin{keybox}[Architecture-Independence Lemma]
The generation law depends \textbf{only} on the unembedding step $\mathbf{z} = W_U \mathbf{h}(x)$, not on how $\mathbf{h}(x)$ is computed. Therefore: Transformer, SSM, Hybrid $\to$ \textbf{same law}.
\end{keybox}

\subsection{$\kappa$ from $W_U$ (zero forward passes)}

Extract $\kappa$ directly from the unembedding matrix: normalize rows, compute nearest-neighbor distances among top-1000 tokens. \textbf{No inference needed.}

\subsection{Results}

\begin{resultbox}[Generation Law Results (22 models)]
\textbf{All models} ($n{=}22$): $r(\kappa, \log\mathrm{CE}) = -0.55$, $p = 0.008$\\
\textbf{Excl.\ LFM outlier} ($n{=}21$): $r = -0.70$, $p < 0.001$\\
\textbf{Fixed-$V$ (Pythia + Mamba)} ($n{=}10$): $r = -0.837$, architecture-independent ($F$-test $p = 0.147$)

Transformers and SSMs lie on the \textbf{same regression line}.
\end{resultbox}

\subsection{Why vocabulary size drops out ($\beta_\mathrm{gen} \approx 0$)}

At each position, only $K_\mathrm{eff} \approx 2$--3 tokens effectively compete (despite $V = 32$K--$152$K). The softmax concentrates probability on a few alternatives, so $\log(V{-}1)$ is irrelevant.

\subsection{Why $R^2$ is only 0.49}

$\kappa$ from $W_U$ captures only output geometry, not context quality $\mathbf{h}(x)$. In classification, $\kappa_\mathrm{nearest}$ captures the \emph{full} race ($R^2 = 0.955$). The generation law's lower $R^2$ is \textbf{predicted} by the theory.

% ===================================================================
\section{Causal Evidence}
% ===================================================================

Five tiers establish that this is causal, not just correlation:

\subsection{Tier 1: Frozen do-interventions}

Move one class centroid in embedding space. Nearest-class intervention: $r = 0.899$. Farthest-class (negative control): $r = 0.000$.

\subsection{Tier 2: Orthogonal factorial}

14 focus classes, clean design: Arm A (nearest) $r = 0.899$; Arm B (control) $r = 0.000$; Arm C (farthest) $r = 0.000$.

\subsection{Tier 3: Cross-modal dose-response}

6 unseen architectures (3 text, 3 vision), frozen $\alpha$, zero refit: 6/6 directional PASS, MAE $= 0.010$.

\subsection{Tier 4: Pre-registered confusion matrix prediction}

\begin{resultbox}[Strongest Causal Test]
Fixed $A = 1.054$, $\tau^* = 0.20$, \textbf{zero degrees of freedom}:

\begin{center}
\begin{tabular}{cccc}
\toprule
$\delta$ & $r$ & Sign accuracy & $p$-value \\
\midrule
1.0 & 0.842 & 92.9\% & $3.8 \times 10^{-50}$ \\
2.0 & 0.808 & \textbf{100\%} & $3.7 \times 10^{-43}$ \\
3.0 & 0.776 & \textbf{100\%} & $8.2 \times 10^{-38}$ \\
\bottomrule
\end{tabular}
\end{center}
$n = 182$ class-pair predictions per level. \textbf{100\% sign accuracy} = every predicted direction correct.
\end{resultbox}

\subsection{Tier 5: Top-$m$ competitor sweep}

$m{=}2$ competitors recover the LOAO slope: $\alpha_2 = 1.523$ vs $\alpha_\mathrm{LOAO} = 1.477$ (2.3\% error). Exactly \textbf{two near-tie competitors} drive the constant.

% ===================================================================
\section{Where the Law Fails (Honest Scope)}
% ===================================================================

\begin{warnbox}[Hard Failures]
\begin{enumerate}[nosep]
\item \textbf{Protein LMs} (ESM-2, 7 models): $\alpha = -1.17$, $r = -0.15$. EC classes too heterogeneous.
\item \textbf{Human fMRI} (NSD, $K{=}12$): 1-NN at chance. Raw BOLD lacks SNR.
\item \textbf{Encoder LOAO}: $\mathrm{CV}_\alpha = 0.42$. Encoder $\alpha$ is not family-invariant.
\end{enumerate}
\end{warnbox}

\textbf{Partial failures}: LODO cross-dataset ($r = 0.125$, expected under three-level universality); large-$K$ regime ($r \approx 0.75$ at $K{=}100$, below noise floor).

% ===================================================================
\section{Three-Level Universality Structure}
% ===================================================================

\begin{keybox}[The Hierarchy (Analogy: Equations of State)]
\textbf{Level 1 --- Form universality} (strongest): The logit-linear structure is EVT-derived and confirmed across NLP, vision, audio, and biological cortex. \emph{Like: $PV = nRT$ holds for all ideal gases.}

\textbf{Level 2 --- Constant within signal type}: $\alpha$ is stable within families (NLP decoders CV $= 2.3\%$) but varies across signal types (NLP $\sim$1.5 vs vision/audio $\sim$4--5). \emph{Like: van der Waals $a,b$ differ by substance.}

\textbf{Level 3 --- Intercept is task-specific}: $C$ encodes dataset difficulty, not predictable across tasks. One calibration measurement required. \emph{Like: ground-state energy differs by system.}
\end{keybox}

% ===================================================================
\section{Open Directions}
% ===================================================================

\subsection{Near-term}
\begin{enumerate}[nosep]
\item Expanded generation law validation (25+ models)
\item External replication by independent groups
\item COLM 2026 submission (abstract Mar 26, paper Mar 31)
\end{enumerate}

\subsection{Medium-term (open theory)}
\begin{enumerate}[nosep, start=4]
\item \textbf{Renormalization theory for $\alpha$ across families}: Why $\alpha_\mathrm{NLP} \approx 1.5$ but $\alpha_\mathrm{ViT} \approx 4.5$ with same $\rho$? The answer lies in $\Sigma_W$ structure.
\item \textbf{Non-asymptotic bounds}: Finite-sample Berry--Esseen for Gumbel order statistics.
\item \textbf{Multilabel extension}: Current law fails for multilabel (e.g., GoEmotions).
\end{enumerate}

\subsection{Long-term (paradigm extensions)}
\begin{enumerate}[nosep, start=7]
\item \textbf{Training objective from the law}: Maximize $\kappa_\mathrm{nearest}$ during training.
\item \textbf{Connection to scaling laws}: If $\kappa \sim C^\gamma$, then $q(C) = \sigma(\alpha C^\gamma + C_0)$.
\item \textbf{Universal renormalization}: Derive $\alpha$ from first principles for \emph{any} architecture.
\end{enumerate}

% ===================================================================
\section{Quick Reference: Key Numbers}
% ===================================================================

\subsection{The Law}
\[
\logit(q_\mathrm{norm}) = \alpha\,\kappa_\mathrm{nearest} - \beta\log(K{-}1) + C
\]

\subsection{Primary constants (NLP decoders, 12 arch.)}
$\alpha = 1.477 \pm 0.034$ (CV $= 2.3\%$), $\beta = 0.478$, $R^2 = 0.955$.

\subsection{Validation scorecard}

\begin{center}
\small
\begin{tabular}{lcc}
\toprule
\textbf{Test} & \textbf{Result} & \textbf{Status} \\
\midrule
12-arch LOAO (NLP) & CV $= 0.023$ & \textcolor{ctigreen}{\textbf{PASS}} \\
RWKV boundary & $2.887 \in [2.43, 3.29]$ & \textcolor{ctigreen}{\textbf{PASS}} \\
Blind OOD (SmolLM2) & $r = 0.817$, $p = 0.013$ & \textcolor{ctigreen}{\textbf{PASS}} \\
ViT-Large & $R^2 = 0.964$ & \textcolor{ctigreen}{\textbf{PASS}} \\
H8+ holdout ($n{=}77$) & $r = 0.879$, MAE $= 0.077$ & \textcolor{ctigreen}{\textbf{PASS}} \\
Audio (7 models) & $r = 0.898$, $p = 0.006$ & \textcolor{ctigreen}{\textbf{PASS}} \\
Bio (32 mice) & 30/32, mean $r = 0.736$ & \textcolor{ctigreen}{\textbf{PASS}} \\
Confusion causal & $r = 0.842$, 100\% sign acc & \textcolor{ctigreen}{\textbf{PASS}} \\
Generation (22 models) & $r = -0.55$ / $-0.70$ & \textcolor{ctigreen}{\textbf{PASS}} \\
Cross-modal $\rho$ & CV $= 1.0\%$ & \textcolor{ctigreen}{\textbf{PASS}} \\
\midrule
Protein ESM-2 & $r = -0.15$ & \textcolor{ctired}{\textbf{FAIL}} \\
Human fMRI & $r = 0.12$, $p = 0.18$ & \textcolor{ctired}{\textbf{FAIL}} \\
Encoder LOAO & CV $= 0.42$ & \textcolor{ctired}{\textbf{FAIL}} \\
\bottomrule
\end{tabular}
\end{center}

\subsection{Zero-parameter predictions}
\begin{itemize}[nosep]
\item $\alpha$ from $\rho$: 1.540 vs 1.477 (+4.3\% error)
\item Confusion matrix: 100\% sign accuracy, zero refit
\item LOAO slope from $m{=}2$: $\alpha_2 = 1.523$ vs $1.477$ (2.3\% error)
\item $K_\mathrm{eff}$ from softmax: $\sim$2--3 ($\beta_\mathrm{gen} \approx 0$)
\end{itemize}

% ===================================================================
\section*{Appendix: The Complete Derivation in One Page}
% ===================================================================

\begin{derivbox}[Seven Steps: Geometry $\to$ Law]
\textbf{Start}: $K$ classes in $\mathbb{R}^d$ with centroids $\{\boldsymbol{\mu}_k\}$ and shared within-class std $\sigma_W$.

\textbf{Step 1} --- Define $\kappa$:
\[
\kappa_\mathrm{nearest}(k) = \min_{j\neq k} \frac{\|\boldsymbol{\mu}_j - \boldsymbol{\mu}_k\|}{\sigma_W\sqrt{d}}
\]

\textbf{Step 2} --- 1-NN as a race:
\[
P(\text{correct} \mid k) = P(M_+ < \min_{j\neq k} M_j)
\]

\textbf{Step 3} --- EVT: $M_+, M_j \to \mathrm{Gumbel}(\mu, s)$ with shared scale $s$.

\textbf{Step 4} --- Gumbel race identity (\textbf{exact}):
\[
P(M_+ < \min_j M_j) = \left[1 + \textstyle\sum_{j\neq k} \exp(-\Delta_{k,j}/s)\right]^{-1}
\]

\textbf{Step 5} --- Nearest-class dominance + sparse hierarchy:
\[
\textstyle\sum \approx (K{-}1)^\beta \exp(-\Delta_\mathrm{nearest}/s), \quad \beta \approx 0.5
\]

\textbf{Step 6} --- Linear gap-to-$\kappa$ map: $\Delta_\mathrm{nearest}/s \approx a_1 \kappa + a_0$.

\textbf{Step 7} --- Combine:
\[
\boxed{\logit(q_\mathrm{norm}) = \alpha\,\kappa_\mathrm{nearest} - \beta\log(K{-}1) + C}
\]

From Euclidean geometry $\to$ EVT $\to$ Gumbel race identity $\to$ logit law. \textbf{QED} (conditional on A1--A3).
\end{derivbox}

\end{document}
